\chapter{Complexity — Strict Circuit}
%%% generated by the blueprint pipeline from TCSlib.Complexity.CircuitComplexity.StrictCircuit
\section{Overview}

Declarations of the Lean module \texttt{TCSlib.Complexity.CircuitComplexity.StrictCircuit}.

\section{Declarations}

\begin{definition}[Strict Gate Condition]
\lean{BoolCircuit.DAGGate.Strict}
\label{BoolCircuit.DAGGate.Strict}
\leanok
\uses{BoolCircuit.DAGGate}
A gate is strict when its list of input vertices has no repetitions and contains exactly
one vertex if the gate is a negation, or exactly two vertices if it is a conjunction or
disjunction.
\end{definition}

\begin{theorem}[Strict gates have fan-in at most two]
\lean{BoolCircuit.DAGGate.Strict.faninTwo}
\label{BoolCircuit.DAGGate.Strict.faninTwo}
\leanok
\uses{BoolCircuit.DAGGate, BoolCircuit.DAGGate.FaninTwo, BoolCircuit.DAGGate.Strict}
\difficulty{4}
Let $g$ be a Boolean gate whose argument indices are pairwise distinct and whose
argument list has length one if $g$ is a NOT gate and length two otherwise. Then its
argument indices are pairwise distinct, a NOT gate has exactly one argument, and $g$ has
at most two arguments.
\end{theorem}

\begin{theorem}[Strict gates have nonempty arguments]
\lean{BoolCircuit.DAGGate.Strict.args_ne_nil}
\label{BoolCircuit.DAGGate.Strict.args_ne_nil}
\leanok
\uses{BoolCircuit.DAGGate, BoolCircuit.DAGGate.Strict}
\difficulty{4}
If a Boolean gate has pairwise distinct argument indices and has exactly one argument
when it is a NOT gate and exactly two otherwise, then its list of argument indices is
nonempty.
\end{theorem}

\begin{definition}[Strict Circuit Condition]
\lean{BoolCircuit.DAGCircuit.IsStrict}
\label{BoolCircuit.DAGCircuit.IsStrict}
\leanok
\uses{BoolCircuit.DAGCircuit, BoolCircuit.DAGCircuit.fanout, BoolCircuit.DAGCircuit.size, BoolCircuit.DAGGate.Strict}
A finite acyclic circuit with a designated output vertex is strict if every gate is
strict and every vertex of index less than the circuit size, other than the output, has
positive fanout.
\end{definition}

\begin{theorem}[Strict circuits have fan-in at most two]
\lean{BoolCircuit.DAGCircuit.IsStrict.isFaninTwo}
\label{BoolCircuit.DAGCircuit.IsStrict.isFaninTwo}
\leanok
\uses{BoolCircuit.DAGCircuit, BoolCircuit.DAGCircuit.IsFaninTwo, BoolCircuit.DAGCircuit.IsStrict, BoolCircuit.DAGGate.Strict.faninTwo}
\difficulty{1}
Let $C$ be a finite Boolean circuit in which every gate has distinct arguments and has
exactly one argument if it is a NOT gate, and exactly two otherwise; suppose every
vertex other than the output has positive fan-out. Then every gate has distinct
arguments, every NOT gate has exactly one argument, and every gate has at most two
arguments.
\end{theorem}

\begin{theorem}[Strict gates have arguments]
\lean{BoolCircuit.DAGCircuit.IsStrict.args_ne_nil}
\label{BoolCircuit.DAGCircuit.IsStrict.args_ne_nil}
\leanok
\uses{BoolCircuit.DAGCircuit, BoolCircuit.DAGCircuit.IsStrict, BoolCircuit.DAGGate.Strict.args_ne_nil}
\difficulty{1}
Let $C$ be a finite Boolean circuit in which every gate has distinct arguments and has
exactly one argument if it is a NOT gate, and exactly two arguments otherwise. Suppose
every vertex other than the output has positive fan-out. Then every gate has a nonempty
argument list.
\end{theorem}

\begin{theorem}[Zero fanout at the circuit boundary]
\lean{BoolCircuit.fanoutIn_last_eq_zero}
\label{BoolCircuit.fanoutIn_last_eq_zero}
\leanok
\uses{BoolCircuit.DAGGate, BoolCircuit.GatesAcyclic, BoolCircuit.fanoutIn}
\difficulty{4}
Let $n$ be a nonnegative integer and let $gs$ be a list of Boolean gates in which every
gate at position $i$ reads only vertices with index less than $n+i$. If $n+|gs|\le v+1$,
then the fanout of vertex $v$ within $gs$ is zero.
\end{theorem}

\begin{theorem}[Output is the last vertex]
\lean{BoolCircuit.DAGCircuit.IsStrict.output_eq}
\label{BoolCircuit.DAGCircuit.IsStrict.output_eq}
\leanok
\uses{BoolCircuit.DAGCircuit, BoolCircuit.DAGCircuit.IsStrict, BoolCircuit.DAGCircuit.fanout, BoolCircuit.DAGCircuit.size, BoolCircuit.fanoutIn_last_eq_zero}
\difficulty{5}
Let $C$ be a finite Boolean circuit in which every gate has distinct arguments and has
exactly one argument if it is a NOT gate, and exactly two otherwise; suppose every
vertex other than the output has positive fan-out. Then the output vertex is the last
vertex, with index $\text{size}(C)-1$.
\end{theorem}

\begin{theorem}[Zero fan-out characterizes the output]
\lean{BoolCircuit.DAGCircuit.IsStrict.fanout_eq_zero_iff}
\label{BoolCircuit.DAGCircuit.IsStrict.fanout_eq_zero_iff}
\leanok
\uses{BoolCircuit.DAGCircuit, BoolCircuit.DAGCircuit.IsStrict, BoolCircuit.DAGCircuit.IsStrict.output_eq, BoolCircuit.DAGCircuit.fanout, BoolCircuit.DAGCircuit.size, BoolCircuit.fanoutIn_last_eq_zero}
\difficulty{4}
Let $C$ be a finite Boolean circuit in which every gate has distinct arguments and has
exactly one argument if it is a NOT gate, and exactly two otherwise. Suppose every
vertex other than the output has positive fan-out. For every vertex $v$ below the size
of $C$, $v$ has fan-out zero if and only if $v$ is the output.
\end{theorem}

\begin{theorem}[No strict circuit with zero inputs]
\lean{BoolCircuit.DAGCircuit.not_isStrict_zero}
\label{BoolCircuit.DAGCircuit.not_isStrict_zero}
\leanok
\uses{BoolCircuit.DAGCircuit, BoolCircuit.DAGCircuit.IsStrict, BoolCircuit.DAGGate.Strict.args_ne_nil}
\difficulty{4}
There is no directed acyclic circuit with zero input vertices in which every gate has
distinct inputs and has exactly one input if it is a NOT gate and exactly two inputs
otherwise, while every vertex other than the output is read by at least one gate.
\end{theorem}

\begin{definition}[Shifted natural index]
\lean{BoolCircuit.shiftPrefix}
\label{BoolCircuit.shiftPrefix}
\leanok
For natural numbers $n$ and $v$, define the shifted index to be $v$ if $v<n$, and $v+3$
otherwise.
\end{definition}

\begin{theorem}[Injectivity of the prefix shift]
\lean{BoolCircuit.shiftPrefix_injective}
\label{BoolCircuit.shiftPrefix_injective}
\leanok
\uses{BoolCircuit.shiftPrefix}
\difficulty{3}
For $n\in\mathbb{N}$, define $f_n(v)=v$ when $v<n$ and $f_n(v)=v+3$ when $v\ge n$. The
map $f_n$ on the natural numbers is injective.
\end{theorem}

\begin{theorem}[Avoiding the shifted prefix]
\lean{BoolCircuit.shiftPrefix_ne}
\label{BoolCircuit.shiftPrefix_ne}
\leanok
\uses{BoolCircuit.shiftPrefix}
\difficulty{3}
Let $n,v,k\in\mathbb{N}$ with $k<3$, and define $f_n(v)=v$ when $v<n$ and $f_n(v)=v+3$
when $v\ge n$. Then $f_n(v)\ne n+k$.
\end{theorem}

\begin{theorem}[Avoiding the first shifted position]
\lean{BoolCircuit.shiftPrefix_ne_one}
\label{BoolCircuit.shiftPrefix_ne_one}
\leanok
\uses{BoolCircuit.shiftPrefix, BoolCircuit.shiftPrefix_ne}
\difficulty{2}
Let $n,v\in\mathbb{N}$, and define $f_n(v)=v$ when $v<n$ and $f_n(v)=v+3$ when $v\ge n$.
Then $f_n(v)\ne n+1$.
\end{theorem}

\begin{theorem}[Shifted prefix avoids the intervening value]
\lean{BoolCircuit.shiftPrefix_ne_two}
\label{BoolCircuit.shiftPrefix_ne_two}
\leanok
\uses{BoolCircuit.shiftPrefix, BoolCircuit.shiftPrefix_ne}
\difficulty{2}
For natural numbers $n$ and $v$, define $s_n(v)=v$ if $v<n$ and $s_n(v)=v+3$ otherwise.
Then $s_n(v)\ne n+2$.
\end{theorem}

\begin{theorem}[Upper bound for the prefix shift]
\lean{BoolCircuit.shiftPrefix_lt}
\label{BoolCircuit.shiftPrefix_lt}
\leanok
\uses{BoolCircuit.shiftPrefix}
\difficulty{3}
Let $n,v,i\in\mathbb{N}$ and suppose $v<n+i$. Define $f_n(v)=v$ when $v<n$ and
$f_n(v)=v+3$ when $v\ge n$. Then $f_n(v)<n+3+i$.
\end{theorem}

\begin{definition}[Constant Circuit Prefix]
\lean{BoolCircuit.constPrefix}
\label{BoolCircuit.constPrefix}
\leanok
\uses{BoolCircuit.DAGGate}
For each natural number $n$, define a three-gate prefix consisting, in order, of a
negation gate reading vertex $0$, a disjunction gate reading vertices $0$ and $n$, and a
conjunction gate reading vertices $0$ and $n$.
\end{definition}

\begin{definition}[Strict Gate Transformation]
\lean{BoolCircuit.strictGate}
\label{BoolCircuit.strictGate}
\leanok
\uses{BoolCircuit.DAGGate, BoolCircuit.shiftPrefix}
For each natural number $n$, define a transformation of gates as follows: an empty-input
conjunction becomes a disjunction reading vertices $n+1,n+2$, and an empty-input
disjunction becomes a conjunction reading those vertices. A one-input conjunction
reading $a$ becomes a conjunction reading the shifted index of $a$ and $n+1$; a
one-input disjunction reading $a$ becomes a disjunction reading the shifted index of $a$
and $n+2$. In all other cases, retain the gate operation and shift every input index by
the prefix shift.
\end{definition}

\begin{theorem}[Arguments of a strict gate]
\lean{BoolCircuit.strictGate_args}
\label{BoolCircuit.strictGate_args}
\leanok
\uses{BoolCircuit.DAGGate, BoolCircuit.shiftPrefix, BoolCircuit.strictGate}
\difficulty{5}
Let $n$ be a natural number and let $g$ be a gate with a natural-number list of input
vertices. Every input vertex of the gate obtained by strictifying $g$ is either the
shifted image of an input vertex of $g$, or one of the two vertices $n+1$ and $n+2$.
\end{theorem}

\begin{theorem}[Strictness of a strictified gate]
\lean{BoolCircuit.strictGate_strict}
\label{BoolCircuit.strictGate_strict}
\leanok
\uses{BoolCircuit.DAGGate, BoolCircuit.DAGGate.FaninTwo, BoolCircuit.DAGGate.Strict, BoolCircuit.shiftPrefix_injective, BoolCircuit.shiftPrefix_ne_one, BoolCircuit.shiftPrefix_ne_two, BoolCircuit.strictGate}
\difficulty{5}
Let $n$ be a natural number and let $g$ be a gate with at most two inputs, no repeated
input vertices, and exactly one input if it is a NOT gate. The gate obtained by
strictifying $g$ has no repeated input vertices and has exactly one input if it is a NOT
gate, and exactly two inputs otherwise.
\end{theorem}

\begin{theorem}[Evaluation of a strict gate]
\lean{BoolCircuit.strictGate_eval}
\label{BoolCircuit.strictGate_eval}
\leanok
\uses{BoolCircuit.DAGGate, BoolCircuit.DAGGate.FaninTwo, BoolCircuit.DAGGate.eval, BoolCircuit.shiftPrefix, BoolCircuit.strictGate}
\difficulty{5}
Let $n$ be a natural number and let $g$ be a gate with at most two inputs, no repeated
input vertices, and exactly one input if it is a NOT gate. For Boolean lists $v$ and
$v'$, suppose the value of $v$ at index $n+1$ (using false for indices outside the list)
is true and its value at index $n+2$ is false. Suppose also that for every input vertex
$a$ of $g$, the value of $v$ at the shifted index of $a$ equals the value of $v'$ at
index $a$, again using false outside either list. Then the strictification of $g$
evaluated on $v$ has the same value as $g$ evaluated on $v'$.
\end{theorem}

\begin{theorem}[Acyclicity after removing constants]
\lean{BoolCircuit.deconst_acyclic}
\label{BoolCircuit.deconst_acyclic}
\leanok
\uses{BoolCircuit.DAGGate, BoolCircuit.GatesAcyclic, BoolCircuit.constPrefix, BoolCircuit.gatesAcyclic_append, BoolCircuit.gatesAcyclic_cons, BoolCircuit.gatesAcyclic_nil, BoolCircuit.shiftPrefix_lt, BoolCircuit.strictGate, BoolCircuit.strictGate_args}
\difficulty{5}
Let $n$ be a positive natural number, and let $G$ be an acyclic list of gates relative
to $n$ inputs. Prefixing the list of three constant-construction gates and then
replacing each gate of $G$ by its strict version produces an acyclic gate list relative
to the same $n$ inputs.
\end{theorem}

\begin{definition}[Circuit Deconstuction]
\lean{BoolCircuit.DAGCircuit.deconst}
\label{BoolCircuit.DAGCircuit.deconst}
\leanok
\uses{BoolCircuit.DAGCircuit, BoolCircuit.constPrefix, BoolCircuit.deconst_acyclic, BoolCircuit.shiftPrefix, BoolCircuit.shiftPrefix_lt, BoolCircuit.strictGate}
Let $C$ be a finite acyclic circuit with $n>0$ input vertices and a valid designated
output. Form a circuit with the same number of inputs by prepending the three-gate
constant prefix, then applying the strict gate transformation with parameter $n$ to each
gate of $C$ in order, and taking the prefix-shifted index of $C$’s output as the new
output.
\end{definition}

\begin{theorem}[Size increase from eliminating constants]
\lean{BoolCircuit.DAGCircuit.deconst_size}
\label{BoolCircuit.DAGCircuit.deconst_size}
\leanok
\uses{BoolCircuit.DAGCircuit.deconst, BoolCircuit.DAGCircuit.size, BoolCircuit.constPrefix}
\difficulty{3}
Let $C$ be a finite acyclic Boolean circuit with $n>0$ inputs. If the size of a circuit
is its number of input vertices plus its number of gates, then eliminating constants
increases its size by exactly three.
\end{theorem}

\begin{theorem}[Strictness of the transformed circuit]
\lean{BoolCircuit.DAGCircuit.deconst_strict}
\label{BoolCircuit.DAGCircuit.deconst_strict}
\leanok
\uses{BoolCircuit.DAGCircuit.IsFaninTwo, BoolCircuit.DAGCircuit.deconst, BoolCircuit.DAGGate.Strict, BoolCircuit.constPrefix, BoolCircuit.strictGate_strict}
\difficulty{4}
Let $n>0$, and let $C$ be a directed acyclic Boolean circuit on $n$ input vertices in
which every gate reads at most two vertices and the gate arguments are distinct, with
every NOT gate reading exactly one vertex. Then every gate in the circuit obtained by
adjoining the three-gate constant prefix and transforming the gates has distinct
arguments and has exactly one argument if it is a NOT gate and exactly two arguments
otherwise.
\end{theorem}

\begin{theorem}[Value preservation under constant removal]
\lean{BoolCircuit.DAGCircuit.deconst_values}
\label{BoolCircuit.DAGCircuit.deconst_values}
\leanok
\uses{BoolCircuit.DAGCircuit.IsFaninTwo, BoolCircuit.DAGCircuit.deconst, BoolCircuit.DAGCircuit.values, BoolCircuit.DAGCircuit.values_getD_gate, BoolCircuit.DAGCircuit.values_getD_input, BoolCircuit.DAGGate.eval, BoolCircuit.DAGGate.eval_pair, BoolCircuit.constPrefix, BoolCircuit.shiftPrefix, BoolCircuit.strictGate, BoolCircuit.strictGate_eval}
\difficulty{6}
Let $n>0$, let $C$ be a directed acyclic Boolean circuit on $n$ input vertices whose
gates have distinct arguments, whose NOT gates each read exactly one vertex, and whose
gates each read at most two vertices. For every input assignment
$x:\{0,\ldots,n-1\}\to\{\mathrm{false},\mathrm{true}\}$ and every vertex $v<n+|C|$, the
value at the shifted vertex in the circuit obtained by adjoining the three-gate constant
prefix and transforming the gates equals the value at $v$ in $C$; values at indices
beyond the available list are interpreted as false.
\end{theorem}

\begin{theorem}[Evaluation preserved by eliminating constants]
\lean{BoolCircuit.DAGCircuit.deconst_eval}
\label{BoolCircuit.DAGCircuit.deconst_eval}
\leanok
\uses{BoolCircuit.DAGCircuit.IsFaninTwo, BoolCircuit.DAGCircuit.deconst, BoolCircuit.DAGCircuit.deconst_values, BoolCircuit.DAGCircuit.eval}
\difficulty{1}
Let $C$ be a finite acyclic Boolean circuit with $n>0$ inputs. Suppose every gate is
well formed, every NOT gate has exactly one argument, and every gate has at most two
arguments. For every Boolean input assignment, the circuit obtained by eliminating
constants has the same output as $C$.
\end{theorem}
